All workstreams share the definitions of \cref{sec:framework} and a common set of conventions.
\begin{itemize}
  \item \emph{Side price and outcome.} Every observation is a side bought at price $\q$, with
  outcome $\won$. Buying YES at YES price $p$ gives $\q=p$; buying NO gives $\q=1-p$.
  \item \emph{Price buckets.} We use twelve buckets of $\q$: $1$--$2$, $3$--$5$, $6$--$10$,
  $11$--$20$, $21$--$35$, $36$--$50$, $51$--$65$, $66$--$80$, $81$--$90$, $91$--$95$, $96$--$98$
  and $99\cents$. Longshots (\LS) are $\q\le10\cents$ and favourites (\FAV) $\q\ge90\cents$.
  \item \emph{Returns.} Returns per contract are $\won-\q-\fee$, and returns per dollar
  follow~\eqref{eq:ret}.
  \item \emph{Fees.} Unless stated otherwise, fees are Kalshi's exact schedule: $0.07$ times the
  series multiplier times $\q(1-\q)$ for takers, and $0.0175\,\q(1-\q)$ for makers on series that
  charge makers. Kalshi computes fees per fill to \$0.000001 and rebates rounding over-payment
  \citep{kalshirounding}, so exact fees are the fees paid.
  \item \emph{Categories.} Kalshi series are grouped into Sports, Weather, Politics (including
  elections and world affairs), Entertainment, Mentions, Science \& Tech, and Other. Polymarket
  markets are mapped to the same groups from their tags. Financial contracts (crypto prices,
  equity indices, interest rates, economic releases and company metrics) are outside the scope
  of this study and are excluded throughout.
  \item \emph{Timing.} Prices used for calibration are observed before the outcome is known:
  \begin{itemize}
    \item at fixed leads before the scheduled resolution;
    \item before the start for sports, except where in-play prices are studied explicitly;
    \item and never after settlement.
  \end{itemize}
  Markets that closed early because the event occurred are timed by their scheduled, not
  realised, close.
  \item \emph{Inference.} Outcomes within an event, such as the legs of a ladder or the teams in
  a game, are dependent. Standard errors are therefore clustered by event
  \citep{cameronmiller2015}, and we report the number of clusters.
  \item \emph{Robustness.} Every main result is shown both market-weighted and volume-weighted
  where the two differ, and separately for the two halves of the sample.
\end{itemize}
